\section{Finite projections and transport of the equation kernel}
\label{sec:finite-rank-kernel-transport}

Free separation does not by itself identify the interacting null space.
We prove the finite-dimensional endpoint, keeping the actual marked
correlation estimates and local equation estimates as explicit inputs.
All spaces below are finite-dimensional over $\mathbb C$.

\subsection{Uniform projections and a common extraction}

Let $V$ have a norm, an isometric conjugate-linear involution $J$, and a
representation $\rho$ of a finite group $\Gamma$ commuting with $J$.
Fix an isomorphism $Q:\mathbb C^m\to V$ from Euclidean space. For any
subspace $U\subset V$, let $\Pi_U$ be the orthogonal projection onto
$Q^{-1}U$. Then
\[
 A_U^0=I-Q\Pi_UQ^{-1},\qquad
 \|A_U^0\|\le C_0:=\|I\|+\|Q\|\,\|Q^{-1}\|
\]
kills $U$ and satisfies $x-A_U^0x\in U$. Any linear map $A$ having
these two properties satisfies $A^2=A$ and $\ker A=U$: apply $A$ to
$x-Ax$ for the first assertion, and use $x=x-Ax$ when $Ax=0$ for
the second.

When $U$ is invariant under $\rho(\Gamma)$ and under $J$, put
\[
 B_U=\frac1{|\Gamma|}\sum_{g\in\Gamma}\rho(g)A_U^0\rho(g)^{-1},
 \qquad A_U=\tfrac12(B_U+JB_UJ).
\]
Each average still kills $U$ and differs from the identity by a map
into $U$. It is therefore an idempotent with kernel $U$. Reindexing
by left multiplication proves that $B_U$ commutes with $\rho(g)$;
commutation of $J$ and $\rho(g)$ proves the same for $A_U$, and
$JA_UJ=A_U$. The composition $JB_UJ$ is complex-linear. Moreover
\[
 \|A_U\|\le\frac{C_0}{|\Gamma|}
       \sum_{g\in\Gamma}\|\rho(g)\|\,\|\rho(g)^{-1}\|.
\]
This bound is independent of $U$. If invariance is known only at late
cutoffs, use $A_U^0$ at earlier ones and take the maximum of the two
bounds. No isometry of the given signed representation is needed.

For every grade $d$, apply this construction to its coefficient space
$V_d$ and finite equation spaces $U_{d,N,p}$. The bounds may depend
on $d$, but not on mesh cutoff $N$ or period index $p$. There is a
common refinement with mesh limits for every fixed $(d,p)$ and then
period limits for every $d$. Indeed, enumerate the countably many
pairs $(d,p)$, take successive subsequences in their compact matrix
balls, and take a diagonal subsequence. Choose its $j$th term from the $j$th nested subsequence,
with index beyond the $j$th index of each of the first $j$ subsequences
and beyond the preceding chosen index. Its tail then belongs to every
fixed earlier subsequence, so every fixed pair retains its limit. Repeat in the period variable for the
countably many grades. Bounded lifted vectors indexed by a comparison
scale and finitely many columns can be included in the same enumeration.
Previously existing limits survive these refinements.

Write $A_{d,p}$ for the inner limits and $T_d$ for the outer limits.
Continuity of multiplication gives $A_{d,p}^2=A_{d,p}$ and $T_d^2=T_d$.
Below the index $c=(N,p)$ follows this ordered filter: first every
sufficiently large selected period, then every sufficiently large
selected mesh cutoff, whose threshold may depend on the period. The
two-stage limits give $A_{d,c}\to T_d$ along this filter. There is no
claim of uniform convergence in $p$.

\subsection{The all-test null space and the local equation input}

Fix a grade, write $V=V_d$, $A_c=A_{d,c}$ and $T=T_d$, and let
$\Theta$ consist of all smooth compactly supported spectator tests
on separated configuration domains with one specified slot of this
grade. Let $f_{\theta,c}\in V^*$ be its finite-cutoff correlation
functional and assume $f_{\theta,c}\to f_\theta$ in operator norm.
Convergence on a fixed basis suffices, since for its dual basis $e_j^*$
\[
 \|f_{\theta,c}-f_\theta\|
 \le\sum_j |(f_{\theta,c}-f_\theta)(e_j)|\,\|e_j^*\|.
\]
Define
\[
                         M=\bigcap_{\theta\in\Theta}\ker f_\theta.
\]
This includes every arity and spectator word, not only the finite pair
measurements used below to bound dimension.

The local equation input is as follows. For every family
$v_c\in\ker A_c$ with subpower inverse-calibrated norm, there is a
comparison insertion $Z_c$ whose separated spectator moments vanish
exactly, and whose difference from the physical insertion of $v_c$
has centered discrepancy price tending to zero on this same filter.
Subpower means that for every $\epsilon>0$ the inverse-calibrated norm
is eventually at most $\exp(\epsilon s_N)$, with $s_N\to\infty$.
The insertions are measurable and bounded at each cutoff, and their
products with the fixed spectator word are integrable.

Section~\ref{ex-column-subpower} proves that bounded coefficient
families have subpower inverse-calibrated norm. The common-cell planting
argument of Section~\ref{subsec:equation-spectator-price} first converts
the root discrepancy into a uniform separated spectator estimate. The
smeared-discrepancy estimate then sends the tested difference from $Z_c$
to zero. The compact
test support is eventually separated on the lattice, so the tested
sum for $Z_c$ is zero. Subtraction of the finite-cutoff integrals is
legitimate by integrability. Consequently
\begin{equation}\label{eq:finite-rank-silent}
 v_c\in\ker A_c,\quad (v_c)\text{ bounded}
       \quad\Longrightarrow\quad f_{\theta,c}(v_c)\longrightarrow0.
\end{equation}
For $v\in\ker T$, the family $v_c=(I-A_c)v$ is bounded, is killed by
$A_c$, and tends to $v$. Hence \eqref{eq:finite-rank-silent} and norm
convergence of the functionals imply
\begin{equation}\label{eq:finite-rank-upper-kernel}
                             \ker T\subseteq M.
\end{equation}
The local equation input is substantive; it does not follow from free
positivity.

It remains applicable after further refinement. Extend a vector family
from the refined grid of mesh and period indices by zero outside that grid. The kernel
condition is preserved. Subpower bounds on the original ordered filter
hold conditionally on membership in the refined grid: monotonicity of
the two index sequences translates the refined thresholds into original
thresholds. Outside that grid the bound is automatic. Apply the original
local estimate to this extension, then restrict back. Thus no control
of the original family at unused cutoff pairs is required.

\subsection{Purchased measurements and exact dimension}

Let $E$ be the free field space, $r=\dim E$, and $b_1,\ldots,b_r$ a
basis. By Section~\ref{sec:ex-lattice-all-grade-free-detectors}, choose
finitely many free measurements $\ell_a\in E^*$ for which
$Lx=(\ell_a(x))_{a=1}^h$ is injective. The additional interacting
comparison consists of vectors $u_{k,j}\in V$ and linear measurements
$g_{k,a}\in V^*$ satisfying
\begin{equation}\label{eq:finite-rank-purchased-matrix}
 g_{k,a}|_M=0,\qquad
             g_{k,a}(u_{k,j})\longrightarrow\ell_a(b_j).
\end{equation}
Here $k$ tends to infinity after taking the mesh and period limits for
each fixed comparison scale. Neither $u_{k,j}$ nor $g_{k,a}$ is required
to converge as $k\to\infty$.

For the application, $g_{k,a}$ is a generic-point correlation linear in
its first insertion, with its other insertions fixed at that $k$.
It vanishes on $M$ because the corresponding correlation is continuous
on the open generic configuration domain and has zero integral against
every compactly supported smooth test there. A nonzero value could be
rotated by a constant complex phase to have positive real part on a
small neighborhood; a nonnegative smooth bump supported there would
then have nonzero integral. This contradiction proves pointwise
vanishing. It uses no point evaluation at a coordinate tie.

The convergence in \eqref{eq:finite-rank-purchased-matrix} is the paid
comparison input. Its precise usable form, after the cutoff limits, is
\[
 |g_{k,a}(u_{k,j})-\ell_{a,k}(b_j)|\le P_{k,a,j}\varepsilon_k,
 \qquad P_{k,a,j}\varepsilon_k\to0,
 \qquad \ell_{a,k}(b_j)\to\ell_a(b_j).
\]
The factors $P_{k,a,j}$ are coefficient norms of the common free moment
profile at the rescaled separated points. Continuity of that finite
profile bounds them. The lifted vectors for fixed $k$ must be bounded
uniformly in both cutoff indices, so the preceding common extraction
retains them. These requirements are not consequences of the finite
matrix argument.

The matrix $(\ell_a(b_j))$ has rank $r$. Choose an $r$-row minor with
nonzero determinant. Its counterpart in $(g_{k,a}(u_{k,j}))$ is nonzero
for all sufficiently large $k$, by entrywise convergence and continuity
of the determinant. A linear relation between the classes $[u_{k,j}]$
in $V/M$, evaluated by these rows, has every coefficient zero. Thus
\begin{equation}\label{eq:finite-rank-lower-dimension}
                              r\le\dim(V/M).
\end{equation}
For $r=0$ this inequality is automatic and no minor is needed.

Assume also the finite equation count
\begin{equation}\label{eq:finite-rank-equation-count}
           \dim V\le\dim\ker A_c+r
           \quad\text{frequently along the same filter}.
\end{equation}
In the application, a label space $W$, equation ideal $\mathcal E$,
and actual calibrated isomorphisms $L_c:W\to V$ satisfy
$\dim(W/\mathcal E)=r$ and $\ker A_c=L_c\mathcal E$. Rank-nullity
then gives equality in \eqref{eq:finite-rank-equation-count}. The free
detector identifies $W/\mathcal E$ with the free space; it does not
establish this separate identity for the actual finite kernel.

If $A,T$ are idempotents with $\|A-T\|<1$, the map
$\ker A\to\ker T$, $v\mapsto(I-T)v$, is injective. Indeed a vector
in its kernel has $Tv=v$ and $Av=0$, so
$\|v\|=\|(A-T)v\|\le\|A-T\|\|v\|$, forcing $v=0$.
Interchange $A,T$ to obtain equal kernel dimensions. Applying this
eventually to $A_c\to T$, and intersecting with the frequent set in
\eqref{eq:finite-rank-equation-count}, gives
$\dim\ker T\ge\dim V-r$. Combining this with
\eqref{eq:finite-rank-upper-kernel} and
\eqref{eq:finite-rank-lower-dimension} proves
\begin{equation}\label{eq:finite-rank-exact-kernel}
                     M=\ker T,\qquad \dim(V/M)=r.
\end{equation}
The purchased classes are therefore eventually a basis of $V/M$.
This supplies the kernel hypothesis in the normalized compact test
argument of Section~\ref{sec:ope-filtration}, provided its projections
and tests are these same ones. That argument then applies to its
unbounded approximating vectors without asserting convergence of
moving physical subspaces.

\subsection{Exact transport and quotient symmetries}

Suppose $\ker A_c=L_c\mathcal E$ as above, with $L_c$ intertwining
the real involutions and signed actions at all sufficiently late
cutoffs. The subspace $\mathcal E$ is invariant under the signed
actions and the real involution, so $U=\ker A_c=L_c\mathcal E$ is
invariant under $\rho(\Gamma)$ and $J$ at those cutoffs; this is the
hypothesis of the averaged construction above. Hence the selected
$A_c$ commute with the actions on $V$. Their limit $T$ commutes too,
by passage to the limit on each fixed vector. Define
\begin{equation}\label{eq:finite-rank-conjugator}
                   Y_c=TA_c+(I-T)(I-A_c).
\end{equation}
Direct multiplication, using both idempotence identities, gives
\[
 Y_c-I=(2T-I)(A_c-T),\qquad Y_cA_c=TA_c=TY_c.
\]
Thus $Y_c\to I$. Choose one late cutoff satisfying all the symmetry
conditions and $\|Y_c-I\|<1$. Its inverse is the norm-convergent
series $\sum_{n\ge0}(I-Y_c)^n$: multiplication of the partial sums
leaves the remainder $(I-Y_c)^{n+1}$, which tends to zero. In particular
$\|Y_c^{-1}\|\le(1-\|Y_c-I\|)^{-1}$. The intertwining identity and
invertibility give $Y_c(\ker A_c)=\ker T$.

Since $Y_c$ is a real polynomial expression in $T,A_c$, it commutes
with both symmetries. Therefore $\mathcal L=Y_cL_c$ is an isomorphism
with $\mathcal L\mathcal E=M$, and descends to
\begin{equation}\label{eq:finite-rank-quotient-isomorphism}
                         W/\mathcal E\simeq V/M
\end{equation}
intertwining the induced conjugations and signed actions. Only one
sufficiently late $L_c$ is used; convergence of $L_c$ is unnecessary.

If a smooth representative of the limiting correlation table has the
same tested correlations, its physical null space equals $M$.
Physical nullity implies tested nullity by integration, while the
converse follows by the smooth bump argument for every separated
spectator word. Hence \eqref{eq:finite-rank-quotient-isomorphism}
identifies that physical quotient with the free quotient. Construction
of the smooth representative and its derivative hierarchy, the actual
calibrated maps $L_c$, the local equation estimate, and the paid
comparisons remain separate inputs. This section proves the finite
projection, dimension and symmetry transport on their common family
and extraction.
